Geben Sie klassische Kenngroeszen zur Datenanalyse

Bilanzgrößen (Prämie, Reserve, Marktpreis), 
Momente (insbesondere Erwartungswert und Varianz), 
Quantile (inbesondere Median und Quartile), Ausreißer

- Anlagen: Validierung der Marktpreise und der entsprechenden Bilanzwerte
- Verpflichtungen: Validierung der Prämien und Reserven und der entsprechenden Bilanzwerte
- Portfolien: Plausibilisierung der Anteile bei Aufteilung in Gruppen (Geschlecht, Tarifgruppen), und Kenngrößen 
(Erwartungswert, Varianz, Quartile, Ausreißer) von stetigen Verteilungen (Prämie, Versicherungssumme, Reserve)

Typische Methoden sind hier das Momente- und Maximum-
Likelihood-Verfahren. Die Anpassung der Daten – und damit die Eignung des Modells
– sollte danach mit einem Verteilungstest geprüft werden (Kolmogorow-Smirnow-
Test oder Q-Q-Plots).

Zur Schätzung von Korrelationen und Parametern von Copulae sind große Daten-
mengen nötig. Statistisch sinnvoll ist das nur im Versicherungskontext auf der Kapi-
talmarktseite umzusetzen.
Die Korrelation kann dort direkt mit Standardverfahren
geschätzt werden. Die Fehler bei der Schätzung sind aber insbesondere bei kleine-
ren Stichproben erheblich. Für andere Anwendungen in der Versicherung werden
üblicherweise Expertenschätzungen verwendet.


Erlaeutern sie den Best Practise Prozess zur Modellimplementierung für Qualitätssicherungsmaßnahme (Testverfahren und Plausibilisierung)

Konsequente Erstellung von Fachkonzept und Testkonzept durch die Fachabteilung
- Organisatorische Trennung von Personen, die die Modelle implementieren von
denen, die die Modelle testen
- Systemtechnisch wird ein reibungsloser Betrieb durch die Trennung von drei
Umgebungen ermöglicht: Test-, Integrations- und Praxisumgebung

r die Durchführung von Tests eignet sich dann das folgende Vorgehen:
1. Testkonzept erstellen: Bereits bei der Erstellung des Fachkonzepts
sollte ein Testkonzept mit erstellt werden. Hier werden mögliche Fehler in der Imple-
mentierung antizipiert, trennscharfe Testfälle definiert und Kriterien zur Prü-
fung festgelegt.
2. Integritätstest (vs. Unit tests)


Geben Sie Beispiele fuer Hypothesentests via Martingale Tests und Leakage

- Martingal-Test: Sind die stochastisch erzeugten Kapitalmarkt-Pfade für die risi-
koneutrale Bewertung tatsächlich Martingale? Notwendige Eigenschaft: Erwar-
tungswert konstant für alle Zeitschritte.
- Leakage: Ergibt der Erwartungswert der Barwerte der Cash-Flows der Anlagen
auch tatsächlich deren Marktpreis?

Im Zeitverlauf erhalten, dem „Funnel of Uncertainty“. 
Der zeitliche Verlauf des Mittelwerts im Vergleich zum Median und die Verläufe z.B. der 1, 5, 25, .., 99 \% 
Quantile lassen auf unplausible Asymmetrien und Inkonsistenzen im Zeitverlauf schließen.

Beisppiele
Profitabilität eines Teilbestandes: Nach gefallenem Zins wird vermutet, dass
ein Teilbestand von aufgeschobenen Renten nicht mehr profitabel bewirtschaftet
werden kann. Muss der Tarif angepasst werden oder nicht? Betrachtung des Funnel
of Uncertainty zum PVFP des Tarifs und zusätzlich Sensitivitäten zu Bestandsverän-
derung und Zinsschwankung.


Was sind die Anforderungen an Solvency II

Für die Dokumentation eines Modells eignet sich grundsätzlich eine Zweiteilung:
- Ein statischer Teil, in dem alle theoretischen Grundlagen, Annahmen und Einschränkungen 
des Modells sowie die methodische Umsetzung des Modells in
allen Facetten beschrieben werden

- Variable d.h. jährlich zu aktualisierende Teile, in denen Berichte zur Durchführung 
der Kalibrierung und Validierung und der entsprechenden Schlussfolgerungen gesammelt werden.

Am Beispiel von Solvency II wird diese Aufteilung detaillierter betrachtet. Die regu-
latorischen Anforderungen an die Dokumentation gemäß §121 VAG:
(1) Der Aufbau und die Funktionsweise des internen Modells sind zu dokumentieren.
Aus dieser Dokumentation muss hervorgehen, dass die Anforderungen der §§ 115
bis 120 VAG eingehalten werden.
(2) Die Dokumentation enthält eine detaillierte Erläuterung der theoretischen Grund-
lagen, der Annahmen sowie der mathematischen und der empirischen Basis, auf
die sich das interne Modell stützt, und beschreibt alle Konstellationen, in denen das
interne Modell nicht wirksam funktioniert.
(3) Versicherungsunternehmen haben alle größeren Veränderungen an ihrem inter-
nen Modell (§111 Absatz 2) zu dokumentieren.

Daraus leitet sich der Inhalte der Dokumentation (insbesondere für Solvency II) ab:
- (statisch) Theoretische Grundlagen, Annahmen und mathematische / empirische Basis,
- (statisch und variabel) Statistische Qualitätsstandards für Wahrscheinlichkeits-
verteilungsprognosen: Vollständigkeit der betrachteten Risiken, Datenbasis,
statistische Methoden und deren Plausibilität,
- (statisch und variabel) Zuordnen von Gewinnen und Verlusten: Vergleich und
Plausibilisierung zwischen Modell und Realität,
- (statisch und variabel) Validierung des Modells.



Geben Sie Beispiele für Definition von Risiko und deren Eigenschaften

- Risiko ist der Effekt den Unsicherheit auf die Erreichung von Zielen hat. (ISO - 31000).
- Risiko als Matrixkombination von Eintrittswahrscheinlichkeit und Auswirkung
- Risiko als Maß für die Volatilität der Rendite einer Investition in der Finanzwirt-
schaft.

Unterscheiden kann man dabei einen einseitigen und einen zweiseitigen Risikobe-
griff.

Bei der einseitigen Sichtweise besteht Risikomanagement konsequenterweise in
der Vermeidung von Risiko. Ansätze zur wertorientierten Steuerung von Unterneh-
men setzen zwingend einen zweiseitigen Risikobegriff voraus.

Gebräuchlich ist die Unterscheidung hinsichtlich der Ungewissheit über die Zukunft
in Risiko (messbar) und der eigentlichen Unsicherheit (Ambiguität).


Was sind die grundlegenden Konzepte von ERM

- Es wird ein ganzheitlicher Ansatz verfolgt. Risiken werden nicht isoliert be-
trachtet und gesteuert. Vielmehr kommen der Berücksichtigung von Abhän-
gigkeiten zwischen den einzelnen Risikokategorien und den Interaktionen der
betroffenen Einheiten eine wesentliche Bedeutung zu.

- Die Betrachtung von Risiken erfolgt zweiseitig
- Auch nicht quantifizierbare Risiken werden berücksichtigt und dabei nach Mög-
lichkeit qualitativ bewertet.

ISO 31000 wurde von der International Organization for Standardization
2009 veröffentlicht mit dem Anspruch Prinzipien und Leitlinien zu formulieren für
das Risikomanagent

COSO hat im Jahr 2004 ihren Standard für ein internes Kontrollsystem erweitert
und auch S\&P (Standard and Poor's) hat in 2005 die Kategorie ERM zu Ihrem Rating hinzugefügt.


Was sind die Schritte des ERM Control Cycles?

Der ERM Control Cycle umfasst die Identifikation, Klassifikation, 
die Bewertung, die Steuerung und die Überwachung von Risiken
Der aktuarielle Control Cycle ist ein Subprozess dessen und besteht aus:
- Spezifikation des Problems,
- Entwicklung einer Lösung,
- Beobachtung des Verlaufs

Eine besondere Herausforderung sind Emerging Risks: 
Dabei handelt es sich um Risiken, die im Beobachtungszeitraum nicht bestanden
oder sich noch nicht verwirklicht hatten.

E.g. Niedrigzinsphasen in der Lebensversicherung



Nenne Sie Beispiele von Stakeholdern und deren Interessen im ERM

- Eigentümer und Investoren haben ein Mindestinteresse am Erhalt der Or-
ganisation und auch an einer angemessenen Rendite in Bezug auf die einge-
gangenen Risiken.
- Für Gläubiger ist der Erhalt der Zahlüngsfähigkeit wesentlich.
- Kunden wie beispielsweise Versicherungsnehmer erwarten, dass die ein-
gegangenen Vertr̈age erfüllt werden. Je nach zugrundeliegendem Vertrag kann
damit auch ein Interesse am wirtschaftlichen Erfolg der Organisation verbun-
den sein, wenn der Kunde daran partizipiert.

Darüber hinaus gibt es noch die
- Ratingagenturen, Verbraucherschützer, Regulatoren, Regierung und Politik, Arbeitnehmer,
Vermittler, Wirtschaftsprüfern, Geschäftsleitung und zuletzt die allgemeine Öffentlichkeit



Was ist das Risikoprofil, Risikoappetit und Risikotoleranz

- Risikoprofil einer Organisation versteht man die Gesamtheit aller Risi-
ken, die mit ihrer Tätigkeit einhergehen

- Der Risikoappetit einer Organisation ist die Beschreibung von Risikozielen und
Risikogrenzen auf Ebene der Gesamtorganisation (e.g. Solvenzquote, risikoadjustierte Ertragsgröße)

- Die resultierende Gesamtheit von Indikatoren und Metriken in den einzelnen Geschäftseinheiten ergibt die Risikotoleranz.
Beispielsweise könnte vorgegeben sein, dass die Wahrscheinlichkeit für eine Belastung von mehr als
100 Mio EUR durch eine Naturkatastrophe kleiner als 0,5 Prozent sein soll.

Für die operative Steuerung der Risiken auf Ebene der risikonehmenden Einheiten
und deren Mitarbeiter müssen konkrete Risikolimiten vorgegeben werden



Geben Sie die detaillierte Klassifikation von Risiken an

Es gibt 3 allg. Risikokategorien in ERM:

Finanzielle Risiken:
- Marktrisiko (Volatilität)
- Kreditrisiko (Bonität)
- Liquiditätsrisiko (Fungibilität)
- Basisrisiko (Hedging ist nicht ausreichend, e.g. CDS, Interest Swaps, Reinsurance) 

Operationale Risiken:
Unangemessenheit oder dem Versagen von internen Prozessen, Mitarbeitern oder Systemen.

- Rechtsrisiko
- Cyber-Risiko (wird kombiniert mit fehlerhafter Programmierung Systemrisiko genannt)
- Managementrisiken, Betrug, kriminelle Handlungen

Strategische Risiken:
Konsequenz nachteiligen Entwicklung des Unternehmens aufgrund strategischer Geschäftsentscheidungen
- die Geschäftsstrategie unzureichend auf das politische, ökonomische oder tech-
nologische Umfeld ausgerichtet ist (insbesondere Unterlassungsentscheidungen (nicht rechtzeitig handeln))

Modell Risiken (speziell für Insurance):
Unzureichendes oder fehlerhaftes Modellieren der Realität 

Diese Risiken haben verschiedene Eigenschaften:

Quantifizierbar
mögliche Modelle sind
- faktor (GLM)
- analytisch (geschlossene Formen)
- Szenariobasisert
- Simulation (ähnlich wie Szenariobasiert, aber es werden Standardverteilungen zur Generierung der Szenarien verwendet)

Als Gegenbeispiel für schwer quantifizierbare Risiken gilt das Reputationsrisiko. Darüber hinaus Risiken die stark korrellieren
und dessen Parameter schwer schätzbar auf Grund mangelnder Daten sind.

Werden die Risiken auf einzelnen Geschäftseinheiten übertragen sollte darauf geachtet werden dass sie folgende
Grundsätze erfüllen:
- Vollständigkeit
- Überschneidungsfreiheit
- Homogenität
- Transparenz



Wie erfolgt der Prozess der Risikoidentifikation und die Beschreibung dieser?


Es gibt mehrere Methoden zur Risikoidentifikation, welche je nach Kontext einsetzbar sind:

- Identifikation von Prozessrisiken durch Prozess-Eigner im Rahmen des IKS (Internen Kontrollsystem)
- Checklisten oder Risikoregister (standardisierte Listen generischer Risiken) als
Basis für unternehmensspezifische Analyse
- Befragung der Mitarbeiter
- externe Einschätzung (z. B. Berater, Rückversicherer)
- statistische Datenanalysen

Die Beschreibung/Dokumentation der einzelnen Risiken sollte beinhalten

- Allg. Zusammenfassung
- Risikokategorie (Financial, Operational, Systemic, Model)
- Zeitraum, Erfassungsdatum
- Verantwortliche
- Brutto-Netto-Bewertung (mit Bewertungsansaetzen)
- Strategie für Kontrolle 

Unter Brutto Risiko versteht man die marginale Verteilung des entstehendes Schadens/Ereignisses .
Unter Netto Risiko die miteinbeziehung von Faktoren wie Steuern und Überschussbeteiligung.



Nennen Sie mögliche Umgangsstrategien für identifizierbare Risiken

- Risikonahme/Risikoakzeptanz
- Risikovermeidung
- Risikobegrenzung
- Risikoübertragung/Risikominderung
- Diversifikation

Konkrete Beispiele sind:
- Produktdesign, Pricing: Risikonahme/Risikovermeidung; vt. Risiken (e.g. Berufsunfähigkeit neben Pensionsgeschäft aber nicht für psychisches Leiden)
- Kapitalanlageprozess, strategische Asset-Allokation: Risikonahme/Risikovermeidung;
Marktrisiken
- Limitsysteme: Risikobegrenzung; grundsätzlich alle Risikokategorien


Was sind Limitsysteme im ERM?

Definition: Limite sind Instrumente, um die gewählte Risikostrategie unter Berück-
sichtigung der Risikotragfähigkeit umzusetzen.

Limitsysteme sind systematisch aufgebaute und an die individuellen Bedürfnisse
eines Unternehmens angepasste Gruppen von Kenngrößen, die für die Steuerung
von Risiken eingesetzt werden.

Sie dienen insbesondere zwei Aspekten:
- Ergebnissteuerung, Erreichen von Planwerten → Controlling
- Mögliche negative Abweichungen von Planwerten → Risikomanagement


Kategorien
- harten Limiten (Risikolimite) => sofortiger Handelsbedarf, schränkt Entscheidungsspielraum ein
- weichen Limiten (Risikoindikatoren) => Mehr Informativ, möglicher Diskussionsbedarf
- präventive Limite => unantastbar


Die konkrete Implementierung ist wie folgt:
1. Ermittlung des Risikodeckungspotential das freigegebene Risikokapital = Risikoappetit
1.1 (freigegebenes Risikokapital = Risikodeckungspotential - Marge)
2. Risikoappetit auf Risikoklassen alloziert
3. Limits pro Klasse festlegen und korrespondierende Schwellenwerte



Systematisieren Sie das verstech. Risiko und nennen Sie dessen Instanzen nach Geschäft

Das versicherungstechnische Risiko ist das Risiko, dass der tatsächliche Aufwand
für Schäden und Leistungen vom erwarteten Aufwand abweicht.

Es besteht aus
- Zufallsrisiko
- Änderungsrisiko 
- Irrtumsrisiko / Modellrisiko


In der LebensVer. sind typische Bestandteile

Sterblichkeitsrisiko,
- das Langlebigkeitsrisiko,
- Invaliditätsrisiko,
- Stornorisiko (bestehend aus Stornoanstiegs-, Stornorückgangs- und Massenstornorisiko)
- Kostenrisiko,
- Katastrophenrisiko;

Schaden-Unfall-Ver.

- Prämienrisiko
- Reserverisiko
- Stornorisiko,
- Katastrophenrisiko

Mögliche Bewertungsmethoden sind:

1. Solvency-II-Standardformel:
1.1 szenariobasiert mit absoluten/prozentualen Anpassungen der betreffenden Wahrscheinlichkeiten oder Kostensätze (Lebens und Krankenver.)
1.2 faktorbasiert mit unternehmensspezifischen Parametern (Schaden-Unfall-Ver)
1.3 modular 
2. multiplikative Gesamtschadenschwankung, bestehend aus vier stochastischen
Komponenten (zufällige Schwankungen, Katastrophen, Schätz-/Irrtumsrisiko Ba-
sistafel, Trend-/Änderungsrisiko)



Systematisieren Sie das markt/kredit-technische und Konzentrations Risiko

- Zinsänderungsrisiko
- Wechselkurs- oder Währungsrisiko
- Aktien und Immobilienrisiko

und der zugehörigen Sensitivität des Portfolios des Unternehmens

Bewertungsansätze sind:
- Monte-Carlo Simulation auf Basis von Kapitalmarkt modellen (geo. BM, CIR Process) mit ökomischen Szenariogenerator 
- faktorbasiert (z. B. Solvency-II)
- Stresstests
- Durationsanalysen

Das Kreditrisiko lässt sich einteilen in
- Ausfallrisiko
- Migrationsrisiko (Änderung der Bonität)
- Spread-Risiko (Änderung des Marktpreises (meist durch Bonität))

Bewertungsansätze
- faktorbasiert (Solvency II mit Bonität und mod. Duration)
- ratingbasiert
- Reduktionsmodelle (stochastisches Modell)

Unter Konzentrationsrisiken versteht man korrelierende Risiken 
innerhalb des Unternehmens.

Bewertungsansätze
- Solvency-II Standardformel
- indirekte Modellierung mit Copulas
- direkte Modellierung mittels Szenarioanalysen


